Use $A=(0,0),B=(6,0),C=(6,6),D=(0,6)$. Then $Q=(0,3),R=(0,1),S=(5,0),N=(6,3)$ and $MP$ is $x=3$. The two other lines are
\[
QS:\ y=3-\dfrac{3}{5}x,\qquad NR:\ y=1+\dfrac{1}{3}x.
\]
Thus $X=\left(3,\dfrac{6}{5}\right)$, $Y=(3,2)$, and $Z=\left(\dfrac{15}{7},\dfrac{12}{7}\right)$. The vertical bases $XY,QR$ have lengths $\dfrac{4}{5},2$, and their distances from $Z$ are $\dfrac{6}{7},\dfrac{15}{7}$. Therefore
\[
\dfrac{[XYZ]}{[QRZ]}=\dfrac{\dfrac{4}{5}\cdot\dfrac{6}{7}}{2\cdot\dfrac{15}{7}}=\boxed{\dfrac{4}{25}}.
\]
\begin{asy}
texpreamble("\usepackage{amsmath,amssymb}");
size(240);
defaultpen(fontsize(9pt));
pair A=(0,0),B=(6,0),C=(6,6),D=(0,6);draw(A--B--C--D--cycle);label("$A$",A,(-1,-1));label("$B$",B,(1,-1));label("$C$",C,(1,1));label("$D$",D,(-1,1));pair M=(3,0),N=(6,3),P=(3,6),Q=(0,3),R=(0,1),S=(5,0),X=(3,1.2),Y=(3,2),Z=(15./7,12./7);draw(M--P);draw(Q--S);draw(N--R);filldraw(X--Y--Z--cycle,lightblue);filldraw(Q--R--Z--cycle,lightgray);label("$M$",M,(0,-1));label("$N$",N,(1,0));label("$P$",P,(0,1));label("$Q$",Q,(-1,0));label("$R$",R,(-1,0));label("$S$",S,(0,-1));label("$X$",X,(1,-1));label("$Y$",Y,(1,1));label("$Z$",Z,(-1,1));
\end{asy}
